Assume \(\mathcal H_{n,d}\ne\varnothing\).  Fix \(P\in\mathcal M_{n,d,L}\), put
\[
 \mathcal G=\sigma\!\left(A,(f_i)_{i\in\mathcal V_n}\right),
 \qquad E_h=\widehat\theta_h-\theta_n,
 \qquad m_h=\mathbb E_P(E_h\mid\mathcal G).
\tag{1}
\]
The minimizer in \({\rm def:exact-optimal-bandwidth}\) exists.  Indeed, across all possible degrees, the finite-lattice score and certificates in \({\rm def:pc-handle}\) have finitely many bandwidth breakpoints.  On every interval beginning at one of its included breakpoints they are constant, apart from the Taylor term, which is continuous and nondecreasing.  The supremum over \(P\in\mathcal M_{n,d,L}\) is also nondecreasing on that interval, so the interval minimum is attained at its included left endpoint.  There are finitely many such endpoints in the nonempty compact interval \(\mathcal H_{n,d}\).  Thus a global minimum is attained, its smallest minimizer \(h^\dagger\) is well defined, and it is fixed by the class rather than by the realized sample.

By \({\rm thm:explicit-pc-conditional-certificate}\),
\[
 |m_h|\leq\overline b_h(A),
 \qquad
 \mathbb E_P\{(E_h-m_h)^2\mid\mathcal G\}\leq\overline s_h(A)^2.
\tag{2}
\]
Conditional Cauchy--Schwarz and (2) imply
\[
 \mathbb E_P(|E_h|\mid\mathcal G)
 \leq |m_h|+\{\mathbb E_P((E_h-m_h)^2\mid\mathcal G)\}^{1/2}
 \leq\overline b_h(A)+\overline s_h(A).
\tag{3}
\]
Taking expectations, setting \(h=h^\dagger\), taking the class supremum, and applying \({\rm def:frontier-envelope}\) yields
\[
 \sup_{P\in\mathcal M_{n,d,L}}
 \mathbb E_P|\widehat\theta_{h^\dagger}-\theta_n|
 \leq\mathfrak u_{n,d,L}.
\tag{4}
\]

Conditional Chebyshev applied to (2) gives
\[
 P\!\left(|E_h-m_h|>\overline s_h(A)/\sqrt\alpha\mid\mathcal G\right)\leq\alpha,
\tag{5}
\]
with a zero left side when \(\overline s_h(A)=0\).  The triangle inequality and \(|m_h|\leq\overline b_h(A)\) turn (5) into coverage by the untruncated interval.  Moreover, \({\rm ass:response-envelope}\) and \({\rm def:target}\) give \(|\theta_n|\leq B\), so intersection with \([-B,B]\) cannot remove \(\theta_n\).  Integrating and setting \(h=h^\dagger\) proves
\[
 \inf_{P\in\mathcal M_{n,d,L}}P\{\theta_n\in\mathcal C_{h^\dagger}\}\geq1-\alpha.
\tag{6}
\]

Before intersection, \({\rm def:honest-interval}\) has half-length exactly \(\overline b_h(A)+\overline s_h(A)/\sqrt\alpha\).  Intersection only shortens it, and therefore
\[
 \operatorname{length}(\mathcal C_h)
 \leq2\max\{1,\alpha^{-1/2}\}\{\overline b_h(A)+\overline s_h(A)\}.
\tag{7}
\]
Taking expectations and the class supremum at \(h=h^\dagger\), then using \({\rm def:frontier-envelope}\), proves
\[
 \sup_{P\in\mathcal M_{n,d,L}}
 \mathbb E_P[\operatorname{length}(\mathcal C_{h^\dagger})]
 \leq2\max\{1,\alpha^{-1/2}\}\mathfrak u_{n,d,L}.
\tag{8}
\]
The sums in \({\rm def:radii}\) explicitly retain every positive realized degree and every isolate.  Equations (4), (6), and (8) prove the theorem; none compares \(\mathfrak u_{n,d,L}\) with \(\mathfrak r^{\mathrm{sup}}_{n,d,L}\).